\documentclass[12pt,a4paper]{amsart}

\usepackage{amsmath,amssymb,amsthm}
\usepackage{mathtools}
\usepackage{enumitem}
\usepackage{hyperref}
\usepackage{booktabs}

%%%─ Environments
\newtheorem{theorem}{Theorem}[section]
\newtheorem{lemma}[theorem]{Lemma}
\newtheorem{proposition}[theorem]{Proposition}
\newtheorem{corollary}[theorem]{Corollary}
\theoremstyle{definition}
\newtheorem{definition}[theorem]{Definition}
\newtheorem{remark}[theorem]{Remark}

%%%─ Macros
\newcommand{\cA}{\mathcal{A}}
\newcommand{\R}{\mathbb{R}}
\newcommand{\cK}{K_{B_c}}
\newcommand{\cKA}{K_{\cA}}
\newcommand{\Ai}{\mathrm{Ai}}
\newcommand{\PK}{P_K}
\newcommand{\PcK}{\widetilde{P}_K^{(c)}}
\newcommand{\Uc}{\widetilde{U}_c}
\newcommand{\Lc}{\mathcal{L}_c^{\mathrm{resc}}}
\newcommand{\xit}{\tilde{\xi}}
\newcommand{\normh}[2]{\|#1\|_{#2}}
\newcommand{\qa}{\mathfrak{a}}
\newcommand{\SK}{S_K}
\newcommand{\HS}{\mathrm{HS}}

\subjclass[2020]{47B34, 45C05, 34E20, 47G10, 60B20}
\keywords{Prolate spheroidal wave functions, Airy kernel,
  Parseval diagonal, uniform spectral tail, Tonelli domination, BR3}

\begin{document}

\title[Paper XVIII: Airy Universality]{%
  Paper~XVIII: Airy Kernel Universality\\[0.3em]
  \large \textsc{Freeze Version} --- BR3 Closed}

\author{\textit{prolate-gram-coercivity programme}}
\date{May 2026}

\begin{abstract}
This paper proves $\cK(\xit,\xit')\to\cKA(\xit,\xit')$ pointwise
(Main Lemma~\ref{lem:main}), closing BR3.
The argument rests on three pillars:
(i) Airy-integral $L^1$-diagonal control
(Proposition~\ref{prop:kernelbridge});
(ii) norm-resolvent spectral convergence
(Kato--Reed--Simon);
(iii) uniform-in-$c$ exponential spectral tail
(Slepian--ORX).
All intermediate steps are self-contained;
no trace of $\mathrm{Id}_{[0,M]}$,
no HS-norm of the full kernel,
and no joint $(c,K)$ uniformity are used.
\end{abstract}

\maketitle
\tableofcontents
\bigskip\hrule\bigskip

%% ============================================================
\section{Setup}
\label{sec:setup}
%% ============================================================

\begin{definition}[Rescaled coordinate]
\label{def:airycoord}
$\xit=c^{2/3}(x_{n^*}^*-x)$,
$(\Uc f)(\xit)=c^{1/3}f(x_{n^*}^*-c^{-2/3}\xit)$,
$\tilde\psi_j=\Uc\psi_{n^*+j}^{(c)}$,
$\cK(\xit,\xit')=c^{-2/3}K_c(x_{n^*}^*-c^{-2/3}\xit,
x_{n^*}^*-c^{-2/3}\xit')$.
Airy window: $c^{-2/3}$;
turning-point shift $\xit_{n^*+j}^*=-j\pi/(2c^{1/3})$ (XVI Lem.~3.2).
\end{definition}

\begin{definition}[Airy operator]
\label{def:airyop}
$\cA=-d^2/d\xit^2+\xit$ on $L^2(\R_+)$,
$\mathcal{D}(\cA)=H^2\cap H^1_0$;
eigenvalues $0<a_1<a_2<\cdots$
($a_k\sim(3\pi k/2)^{2/3}$);
eigenfunctions $u_k=\Ai(\cdot-a_k)/\normh{\Ai(\cdot-a_k)}{L^2(\R_+)}$,
$\normh{u_k}{L^2(\R_+)}=1$.
\end{definition}

\begin{remark}[Compact embeddings, $n=1$]
\label{rem:compact}
On bounded $[0,M]$:
$H^2\hookrightarrow H^1\hookrightarrow C^0\hookrightarrow L^2$,
all compact (Rellich--Kondrachov);
$\normh{f}{L^\infty}\leq C_M\normh{f}{H^1}$ (Morrey).
\end{remark}

\begin{proposition}[Airy kernel; $L^1$ diagonal]
\label{prop:kernelbridge}
\textbf{(a) Integral representation.}
$\cKA(\xit,\xit')
=\sum_k u_k(\xit)u_k(\xit')
=\int_0^\infty\Ai(\xit+s)\Ai(\xit'+s)\,ds$
(Tracy--Widom~\cite{TracyWidom1994});
$\cKA\in L^1(\R_+^2)$ (Airy decay).

\medskip
\textbf{(b) $L^1$ diagonal via Airy decay.}
For $\xit\geq 0$:
\begin{equation}\label{eq:diagbound}
  \cKA(\xit,\xit)
  =\int_0^\infty\Ai(\xit+s)^2\,ds
  \leq C'(1+\xit)^{-1/2},
\end{equation}
using $|\Ai(t)|\leq C(1+t)^{-1/4}e^{-\frac{2}{3}t^{3/2}}$
(Olver~\cite{Olver1997} Ch.~11).
Hence $\cKA(\xit,\xit)\in L^1([0,M])$ for every $M<\infty$.
\emph{No semigroup, no trace of identity.}

\medskip
\textbf{(c) Parseval identification.}
For a.e.\ $\xit\in\R_+$:
$\cKA(\xit,\xit)=\sum_k u_k(\xit)^2$,
convergence in $L^1([0,M])$
(Bessel + DCT with dominant~\eqref{eq:diagbound}).
\end{proposition}

\begin{lemma}[Main Lemma]\label{lem:main}
For fixed $\xit,\xit'\geq 0$:
$\cK(\xit,\xit')\to\cKA(\xit,\xit')$.
\end{lemma}

%% ============================================================
\section{Formal Poset Structure}
\label{sec:depgraph}
%% ============================================================

\begin{definition}[Internal dependency poset]
\label{def:poset}
Let $\mathcal{N}$ be the set of all internal results of this paper.
Define $A\leq B$ iff $A$ appears in the proof of $B$ within
$\mathcal{N}$.
External results form a separate set $\mathcal{E}$,
entering $\mathcal{N}$ only as labeled inputs.
\end{definition}

\begin{center}
\renewcommand{\arraystretch}{1.2}
\begin{tabular}{l l l l}
\toprule
\textbf{Node} & \textbf{Intro} & \textbf{Internal} & \textbf{From $\mathcal{E}$} \\
\midrule
$d_1$: Def.~\ref{def:airycoord} & $c$ & --- & XVI.3.2 \\
$d_2$: Def.~\ref{def:airyop} & $a_k$ & --- & Airy zeros \\
$r_0$: Rem.~\ref{rem:compact} & $M$ & --- & Sob., Rel.-K. \\
$p_0$: Prop.~\ref{prop:kernelbridge} & --- & $d_2$ & Olver, TW \\
$p_1$: Prop.~\ref{prop:formbounded} & $C_0$ & $d_1$ & XVI, Kato.VI \\
\hline
$\ell_0$: Lem.~\ref{lem:slbound} & $C_K,c_0(K)$ & $p_1$ & Kato.VI.2.21 \\
$\ell_A$: Lem.~\ref{lem:langer} & $\delta_0,c_0(K,M)$ & $\ell_0,r_0$ & Olver.11.3 \\
$\ell_L$: Lem.~\ref{lem:linf} & $C(M,K)$ & $\ell_0,r_0$ & $H^2$-reg. \\
$\ell_{C0}$: Lem.~\ref{lem:slepianunif} & $K_0,c_1,\alpha$ & --- & L-W, Slepian, ORX \\
$p_2$: Prop.~\ref{prop:normresolvent} & $z$ & $p_1$ & RS.VIII.23 \\
\hline
$c_1$: Cor.~\ref{cor:spectral}(i--iii) & & $p_2$ & Kato.IV \\
$\ell_2$: Lem.~\ref{lem:H2} & & $\ell_A,r_0$ & --- \\
$\ell_B$: Lem.~\ref{lem:sturmunif} & & $\ell_0,c_1,\ell_2$ & Zettl.8.3 \\
$\ell_{BT}$: Lem.~\ref{lem:besseltail} & & $p_0$ & RS.II.4 \\
$\ell_D$: Lem.~\ref{lem:domconv} & & $\ell_L,\ell_{C0},\ell_2,p_0,\ell_{BT}$ & DCT \\
$\ell_I$: Lem.~\ref{lem:I2b} & & $\ell_A$ & Airy decay \\
\hline
$m$: Main Lem.~\ref{lem:main} & & $\ell_D,c_1,\ell_I$ & XVII.2.1 \\
$T$: Thm.~\ref{thm:br3} & & $m$ & XVII.4.2 \\
\bottomrule
\end{tabular}
\end{center}

\begin{remark}[Feedback-stable, no cycle]
\label{rem:feedback}
$\ell_B$ uses $c_1$ (from $p_2$, which depends only on $p_1$)
and $\ell_2$ (from $\ell_A$ only).
$p_2\not\leq\ell_B$: no directed cycle in $(\mathcal{N},\leq)$.
\end{remark}

%% ============================================================
\section{Closing Module}
\label{sec:closing}
%% ============================================================

\begin{proposition}[Intertwining]
\label{prop:formbounded}
$\Lc=\cA+c^{-1/3}R_c$,
$|W_c|\leq C(1+|\xit|)^{1/2}$,
$\normh{R_c}{}\leq C_0$;
form-rel.\ bounded, rel.\ bound $0$
(Kato~\cite{Kato1966} VI.1.38); type~(A).
\end{proposition}

\begin{lemma}[Global eigenvalue bound]
\label{lem:slbound}
$a_{j,c}=a_j+O_j(c^{-1/3})$;
$a_{j,c}\leq C_K:=a_K+1$ for $c\geq c_0(K)$.
\end{lemma}
\begin{proof}
$\qa_c=\qa_0+c^{-1/3}\langle W_c\cdot,\cdot\rangle$ on $H^1_0(\R_+)$;
Hardy--Rellich + Kato~\cite{Kato1966} VI.2.21.
\end{proof}

\begin{remark}[Uniformity chain]
\label{rem:uniform}
$K\xmapsto{\ell_0}C_K=a_K{+}1
\xmapsto{}\delta_0(K)=2C_K{+}2
\xmapsto{\ell_A}c_0(K,M)$.
All constants in Lemma~\ref{lem:langer}
are $C(M,K)$, uniform in $j\leq K$ and $c\geq c_0(K,M)$.
\end{remark}

\begin{remark}[Spectral $\neq$ coercivity]
\label{rem:coercep}
Lemma~\ref{lem:slbound}: global $a_{j,c}\leq C_K$.
Lemma~\ref{lem:langer} Step~2: local form inequality
on $[\delta_0,M]$, using $C_K$ as a fixed numerical input only.
\end{remark}

\begin{lemma}[Langer + coercive]
\label{lem:langer}
$\tilde\psi_j=c_j[\Ai(\cdot-\xit_{n^*+j}^*)+E_j^{(c)}]$,
$\sup_{[0,M]}|E_j^{(c)}|\leq C'(M,K)c^{-1/3}$,
$c_j=1+O(C(K)c^{-1/3})$;
uniform in $j\leq K$, $c\geq c_0(M,K)$.
\end{lemma}
\begin{proof}
$\delta_0=2C_K+2$.
\textbf{Step 1 (Olver).}
$\sup_{I_c}|E_j^{(c)}|=O(c^{-1})$ (Olver~\cite{Olver1997} Thm.~11.3).
\textbf{Step 2 (form ineq.).}
$q_c\geq C_K+1=:\mu(K)>0$ on $[\delta_0,M]$;
$\qa_c[f]\geq\mu(K)\normh{f}{H^1}^2$.
\textbf{Step 3 ($H^{-1}$, $n=1$).}
Lax--Milgram: $\sup|e_j|\leq C'(M,K)c^{-1/3}$.
\end{proof}

\begin{remark}[$K$-dependence of $C'(M,K)$: geometric origin]
\label{rem:langerkdep}
The constant $C'(M,K)$ in Lemma~\ref{lem:langer} depends on $K$
through a single geometric chain:
\[
  K \;\xmapsto{\text{Airy asymptotics}}\;
  a_K \sim \Bigl(\tfrac{3\pi K}{2}\Bigr)^{2/3}
  \;\xmapsto{\ell_0}\;
  C_K = a_K + 1
  \;\xmapsto{\text{Step 2}}\;
  \delta_0(K) = 2C_K + 2
  \;\xmapsto{\text{Lax--Milgram}}\;
  C'(M,K).
\]
Concretely: the coercivity threshold $\mu(K) = C_K + 1$ and the
exclusion interval $[0,\delta_0(K))$ both grow like $K^{2/3}$,
which feeds a polynomial $K$-dependence into the Lax--Milgram constant.
A direct tracking of constants in Steps~1--3 shows
\begin{equation}\label{eq:Cprimescale}
  C'(M,K) = O\bigl(M K^{2/3}\bigr) \qquad (K\to\infty,\; M\text{ fixed}),
\end{equation}
where the implicit constant is independent of $c$.

\medskip
\textbf{Consequence for Paper~XIX.}
This polynomial $K$-growth is the \emph{only} source of
$K$-dependence in the Langer-regime rate; it is
\emph{geometrically} driven (Airy zero spacing $a_K\sim K^{2/3}$),
not analytically critical.
In particular:
\begin{itemize}[nosep]
  \item No uniform-in-$K$ Langer estimate is needed for BR3;
    only the fixed-$K$ bound is used (XIX Lemma~\texttt{lem:langer\_rate}).
  \item With $C'(M,K) = O(K^{2/3})$, the XIX constant becomes
    $C_1(M,K) = O(K^{5/3} M)$,
    confirming polynomial (not exponential) growth.
  \item The logarithmic coupling $K_\mathrm{opt}(c)=(\log c)/(3\alpha)$
    gives $C_1(M,K_\mathrm{opt}) = O((\log c)^{5/3})$,
    hence the coupled rate is
    $O(c^{-1/3}(\log c)^{5/3})$ --- still $o(c^{-1/3+\varepsilon})$
    for every $\varepsilon > 0$.
\end{itemize}
\end{remark}

\begin{lemma}[Uniform $L^\infty$]
\label{lem:linf}
$\normh{\tilde\psi_j}{L^\infty([0,M])}\leq C(M,K)$,
uniform in $j\leq K$, $c\geq c_0(M,K)$.
\end{lemma}

\begin{lemma}[Uniform-in-$c$ spectral tail]
\label{lem:slepianunif}
For each $\varepsilon>0$ there exist
$\alpha=\alpha(\varepsilon)>0$, $K_0=K_0(\varepsilon)$,
and $c_1=c_1(\varepsilon)$ with $\alpha$ independent of $c$, such that
\begin{equation}\label{eq:uniftail}
  \lambda_{n_c+k}^{(c)}\leq e^{-\alpha k}
  \quad\text{for all }k\geq K_0,\ c\geq c_1.
\end{equation}
Consequently
$\sup_{c\geq c_1}\sum_{j>K_0}\lambda_{n^*+j}^{(c)}
\leq\frac{e^{-\alpha(K_0+1)}}{1-e^{-\alpha}}<\varepsilon$.
\end{lemma}
\begin{proof}
Write $n_c=\lfloor2c/\pi\rfloor$. Fix $\delta\in(0,\tfrac{1}{3})$.

\textbf{Inner transition} ($K_0\leq k\leq c^{1/3+\delta}$).
ORX~\cite{Osipov2013} \S4.2.1:
$\lambda_{n_c+k}^{(c)}\leq e^{-\beta_1 k}$,
$\beta_1=\beta_1(\delta)>0$ uniform in $c\geq c_1(\delta)$.

\textbf{Deep tail} ($k>c^{1/3+\delta}$).
Slepian~\cite{Slepian1978} Thm.~1:
$\lambda_{n_c+k}^{(c)}<2e^{-\beta_2 k}$,
$\beta_2=\beta_2(\varepsilon)>0$ uniform in $c\geq c_1(\varepsilon)$.

\textbf{Gluing.}
Both regimes cover $k\geq K_0$ without gap;
$\alpha:=\min(\beta_1,\beta_2)>0$
satisfies~\eqref{eq:uniftail} uniformly.
\end{proof}

\begin{remark}[Iterated limit; no joint uniformity needed]
\label{rem:iterated}
For each $\varepsilon>0$:
(i) choose $K_0$ from Lemma~\ref{lem:slepianunif}
(uniform in $c\geq c_1$);
(ii) choose $c_0=c_0(K_0,M)$ from Lemma~\ref{lem:langer}.
For $c\geq\max(c_0,c_1)$: total error $<\varepsilon$.
\emph{No joint uniform convergence in $(c,K)$ is required;
only the sequential limit is used.}
\end{remark}

\begin{lemma}[Sturm ordering]
\label{lem:sturmunif}
$a_{j,c}<a_{j',c}$ ($j<j'$);
$\tilde\psi_j\leftrightarrow u_j$ stable.
\end{lemma}
\begin{proof}
(i) Zettl~\cite{Zettl2005} 8.3.1.
(ii) Gap: $\ell_0$ + Airy spacing.
(iii) $\tilde\psi_j\to u_j$ in $C^0([0,M])$
(Lemma~\ref{lem:H2}; no loop, Remark~\ref{rem:feedback}).
\end{proof}

%% ============================================================
\section{Spectral Convergence}
\label{sec:nrc}
%% ============================================================

\begin{proposition}[Norm-resolvent]
\label{prop:normresolvent}
$(\Lc-z)^{-1}-(\cA-z)^{-1}
=-(\cA-z)^{-1}c^{-1/3}R_c(\Lc-z)^{-1}$;
norm $\leq C_0c^{-1/3}/(\mathrm{Im}\,z)^2$;
$P_{j,c}\to P_j$ strongly (RS~\cite{ReedSimon4} VIII.23--24).
\end{proposition}

\begin{corollary}[Spectral convergence]
\label{cor:spectral}
(i) $|a_{j,c}-a_j|=O_j(c^{-1/3})$;
(ii) $\normh{P_{j,c}-P_j}{}=O_j(c^{-1/3})$;
(iii) $\normh{\PcK-\PK}{}=O_K(c^{-1/3})$;
(iv) $\normh{\tilde\psi_j}{H^1}\leq C_K$.
\end{corollary}

\begin{lemma}[$C^0$ convergence]
\label{lem:H2}
$\normh{\tilde\psi_j-u_j}{C^0([0,M])}\leq C'(M,K)c^{-1/3}$.
\end{lemma}
\begin{proof}
Lemma~\ref{lem:langer}: $c_j\to 1$,
$\xit_{n^*+j}^*\to 0$,
$\sup|E_j^{(c)}|=O(c^{-1/3})$.
\end{proof}

%% ============================================================
%% KEY LEMMA: Bessel tail on restricted interval
%% ============================================================
\begin{lemma}[Bessel tail on $[0,M]$]
\label{lem:besseltail}
Let $\{u_j\}$ be the $L^2(\R_+)$-ONB of Airy eigenfunctions
(Definition~\ref{def:airyop}). For every $M<\infty$:
\begin{equation}\label{eq:besseltailconv}
  \sum_{j>K}\normh{u_j}{L^2([0,M])}^2 \to 0
  \quad\text{as }K\to\infty.
\end{equation}
\end{lemma}
\begin{proof}
Set $b_j:=\normh{u_j}{L^2([0,M])}^2\in[0,1]$.
The partial sums $T_K:=\sum_{j=1}^K b_j$
are monotone non-decreasing and bounded:
\[
  T_K \leq \int_0^M\cKA(\xit,\xit)\,d\xit
  \leq C'\int_0^M(1+\xit)^{-1/2}\,d\xit <\infty
\]
(Bessel's inequality in $L^2(\R_+)$, RS~\cite{ReedSimon2} Prop.~II.4,
plus~\eqref{eq:diagbound}).
Hence $T_K\nearrow T_\infty<\infty$,
and $\sum_{j>K}b_j = T_\infty-T_K\to 0$.
\end{proof}

\begin{lemma}[Spectral sum convergence]
\label{lem:domconv}
For fixed $\xit,\xit'\in[0,M]$:
$\SK(\xit,\xit'):=\sum_{j=1}^K\tilde\psi_j(\xit)\tilde\psi_j(\xit')
\to\cKA(\xit,\xit')$
($c\to\infty$ then $K\to\infty$).
\end{lemma}
\begin{proof}
Fixed $K$: Lemma~\ref{lem:H2}.
$K\to\infty$: Remark~\ref{rem:dct} + Remark~\ref{rem:iterated}.
\end{proof}

\begin{remark}[DCT: Cauchy--Schwarz + Tonelli + Bessel]
\label{rem:dct}
\textbf{Layer 1 (fixed $K$).}
$|\SK|\leq K\cdot C(M,K)^2$ (Lemma~\ref{lem:linf}).

\medskip
\textbf{Layer 2 ($K\to\infty$).}

\textit{Step 1: Cauchy--Schwarz domination.}
For all $K\geq 1$ and $(\xit,\xit')\in[0,M]^2$:
\begin{equation}\label{eq:CStail}
  |\cKA(\xit,\xit')-\SK(\xit,\xit')|
  \leq
  [\cKA(\xit,\xit)]^{1/2}
  [\cKA(\xit',\xit')]^{1/2}
  =: h(\xit,\xit').
\end{equation}

\textit{Step 2: Tonelli factorisation.}
\begin{equation}\label{eq:Tonelli}
  \normh{h}{L^1([0,M]^2)}
  = \Bigl(\int_0^M\cKA(\xit,\xit)^{1/2}\,d\xit\Bigr)^2
  < \infty
\end{equation}
(using~\eqref{eq:diagbound}; $h\notin L^1(\R_+^2)$ globally).

\textit{Step 3: $L^2$-convergence via Lemma~\ref{lem:besseltail}.}
Write $\SK=\sum_{j=1}^K P_j\otimes P_j$,
$P_j=\langle u_j,\cdot\rangle u_j$ in $L^2(\R_+)$.
By Cauchy--Schwarz on the kernel expansion:
\begin{equation}\label{eq:L2conv}
  \normh{\cKA-\SK}{L^2([0,M]^2)}^2
  \leq \sum_{j>K}\normh{u_j}{L^2([0,M])}^2
  \;\to\; 0
\end{equation}
by Lemma~\ref{lem:besseltail}.
Since the $L^2([0,M]^2)$-limit is unique,
a standard subsequence extraction yields
$\SK_n(\xit,\xit')\to\cKA(\xit,\xit')$
a.e.\ on $[0,M]^2$.

\textit{Step 4: Lebesgue DCT.}
Bound~\eqref{eq:CStail} with $h\in L^1([0,M]^2)$
+ a.e.\ convergence $\Rightarrow$ $\SK\to\cKA$ in $L^1([0,M]^2)$.
\emph{No trace of $\mathrm{Id}_{[0,M]}$;
no HS-norm;
no Mercer monotonicity;
no $S_K$-monotonicity in $K$.}
\end{remark}

%% ============================================================
\section{Proof of Main Lemma}
\label{sec:proof}
%% ============================================================

\begin{lemma}[Kernel domination]
\label{lem:I2b}
$|\cK|\leq C(M)(1+\xit)^{-1/4}(1+\xit')^{-1/4}$
on $[0,M]^2$; Airy decay outside.
\end{lemma}

\textbf{Proof of Main Lemma~\ref{lem:main}.}
\textit{Step~1.} Fermi tail: XVII Lem.~2.1 + Lem.~\ref{lem:I2b}.
\textit{Step~2.} Projector: Cor.~\ref{cor:spectral}(i--iii).
\textit{Step~3.} Fixed-$K$: Lem.~\ref{lem:domconv}.
\textit{Step~4.} $K\to\infty$: Lem.~\ref{lem:slepianunif}
+ Rem.~\ref{rem:dct} + Rem.~\ref{rem:iterated}.$\hfill\square$

%% ============================================================
\section{BR3}
\label{sec:br3}
%% ============================================================

\begin{theorem}[BR3]\label{thm:br3}
$\normh{\cK-\cKA}{L^2(w\otimes w)}\to 0$
(Lem.~\ref{lem:main} + XVII Thm.~4.2).
\end{theorem}

\begin{center}
\renewcommand{\arraystretch}{1.3}
\begin{tabular}{p{7.5cm} p{0.9cm} p{2.9cm}}
\toprule
\textbf{Claim} & \textbf{Status} & \textbf{Source} \\
\midrule
$L^1$ diagonal (Airy integral, no semigroup)
  & \checkmark & Prop.~\ref{prop:kernelbridge}(b) \\
Tonelli: $\normh{h}{L^1}$ on $[0,M]^2$ only
  & \checkmark & Eq.~\eqref{eq:Tonelli} \\
Bessel tail $\to 0$ on $[0,M]$ (standalone)
  & \checkmark & Lem.~\ref{lem:besseltail} \\
$L^2([0,M]^2)\Rightarrow$ subseq.\ a.e.
  & \checkmark & Rem.~\ref{rem:dct} Step~3 \\
Slepian: $\alpha=\min(\beta_1,\beta_2)$
  & \checkmark & Lem.~\ref{lem:slepianunif} \\
Iterated limit, no joint $(c,K)$ uniformity
  & \checkmark & Rem.~\ref{rem:iterated} \\
Langer ($c^{-2/3}$ window)
  & \checkmark & Lem.~\ref{lem:langer} \\
$K$-dep.\ of $C'(M,K)$: geometric, $O(K^{2/3})$
  & \checkmark & Rem.~\ref{rem:langerkdep} \\
Poset DAG, no cycle
  & \checkmark & Rem.~\ref{rem:feedback} \\
Main Lemma & \checkmark & Steps 1--4 \\
BR3 & \checkmark & XVII + XVIII \\
\bottomrule
\end{tabular}
\end{center}

\begin{thebibliography}{99}
\bibitem{paper16} \textit{Paper~XVI}, 2026.
\bibitem{paper17} \textit{Paper~XVII}, 2026.
\bibitem{Kato1966} T.~Kato,
  \textit{Perturbation Theory for Linear Operators}, Springer, 1966.
\bibitem{ReedSimon2} M.~Reed, B.~Simon,
  \textit{Methods of Modern Mathematical Physics, Vol.~II},
  Academic Press, 1975.
\bibitem{ReedSimon4} M.~Reed, B.~Simon,
  \textit{Methods of Modern Mathematical Physics, Vol.~IV},
  Academic Press, 1978.
\bibitem{Zettl2005} A.~Zettl,
  \textit{Sturm--Liouville Theory}, AMS, 2005.
\bibitem{Olver1997} F.~Olver,
  \textit{Asymptotics and Special Functions}, AK Peters, 1997.
\bibitem{TracyWidom1994} C.~Tracy, H.~Widom,
  Level-spacing distributions and the Airy kernel,
  \textit{Comm.\ Math.\ Phys.}\ \textbf{159} (1994), 151--174.
\bibitem{LandauWidom1980} H.~Landau, H.~Widom,
  Eigenvalue distribution of time and frequency limiting,
  \textit{J.\ Math.\ Anal.\ Appl.}\ \textbf{77} (1980), 469--481.
\bibitem{Slepian1978} D.~Slepian,
  Prolate spheroidal wave functions, Fourier analysis, and
  uncertainty~V,
  \textit{Bell Syst.\ Tech.\ J.}\ \textbf{57} (1978), 1371--1430.
\bibitem{Osipov2013} A.~Osipov, V.~Rokhlin, H.~Xiao,
  \textit{Prolate Spheroidal Wave Functions of Order Zero},
  Springer, 2013.
\bibitem{Davies1995} E.B.~Davies,
  \textit{Spectral Theory and Differential Operators},
  Cambridge, 1995.
\bibitem{Simon2005} B.~Simon,
  \textit{Trace Ideals and Their Applications}, 2nd ed., AMS, 2005.
\bibitem{AbramowitzStegun1964} M.~Abramowitz, I.~Stegun,
  \textit{Handbook of Mathematical Functions}, NBS, 1964.
\end{thebibliography}

\end{document}
